\documentclass[a4paper]{article}

\usepackage{graphicx}
\usepackage{amsmath,amssymb}
\usepackage{minted}
\usepackage{multirow}
\usepackage{float}
\usepackage{todonotes}
\usepackage{forest}
\usepackage{xstring}
\usepackage{xcolor}
\usepackage[most]{tcolorbox}
\usepackage{xparse}
\usepackage{natbib}

\usetikzlibrary{graphs,graphdrawing}
\usegdlibrary{layered}

% for external graphviz
\input{/home/donkarlo/Dropbox/repo/nd_language_ai_project/src/nd_language_ai/formal/computer/programming/tex/latex/code_clips/external_graphviz.tex}

% for tags
\input{/home/donkarlo/Dropbox/repo/nd_language_ai_project/src/nd_language_ai/formal/computer/programming/tex/latex/part/control_sequence/kind/semtag/semtag.tex}

% for links
\input{/home/donkarlo/Dropbox/repo/nd_language_ai_project/src/nd_language_ai/formal/computer/programming/tex/latex/code_clips/seehere.tex}

\title{Plan}
\date{}

\begin{document}
   \maketitle
   \section{Definition}
      A plan is a sequence of actions intended to achieve a goal from a given
      initial state \cite{mcdermott-1999-using-regression-match-graphs-to-control-search-in-planning}.
      
      \cite{mcdermott-1999-using-regression-match-graphs-to-control-search-in-planning} formulates classical planning as the problem of finding a sequence of actions that achieves a specified goal. A planning problem assumes knowledge about the initial state and the effects of available actions. A plan therefore represents an ordered course of actions leading toward the desired goal state. The proposed method searches through plan prefixes and progressively extends them toward complete plans. This formulation provides a formal basis for representing goal-directed plans in artificial intelligence.

      A plan is typically any list of steps, with details of timing and resources, used to achieve an objective. \seehere{https://en.wikipedia.org/wiki/Plan}
      %\bibliographystyle{plainnat}
      %\bibliography{/home/donkarlo/Dropbox/repo/nd_shared_research/bibliography/bibliography.bib}
\end{document}